\begin{exercice}\label{exo1-4} (\minsyndical)

On effectue une moyenne sur 30 ans des relev\'es de temp\'erature mensuelle en 
$^o$C \`a New-York et \`a San Francisco. On obtient le tableau qui suit:\\

\begin{tabular}{|c|cccccccccccc|}
\hline
Mois & Jan & Fev & Mar & Avr & Mai & Jui & Jul & Aou & Sep & Oct & Nov & Dec \\ \hline
New-York & 0 & 1 & 5 & 12 & 17 & 22 & 25 & 24 & 20 & 14 & 8 & 2 \\ \hline
San Francisco & 9 & 11 & 12 & 13 & 14 & 16 & 17 & 17 & 18  & 16 & 13 & 9 \\ \hline
\end{tabular}\\

Calculer :

\begin{enumerate}

\item La moyenne empirique $\bar{x}_n$.

\item La m\'ediane empirique $\tilde{x}_n$.

\item L' \'ecart-type empirique $s_n$.

\item Le coefficient de variation empirique $s_n / \bar{x}_n$.

%\item Les quartiles $\tilde{q}_{n,1/4},\tilde{q}_{n,1/2},\tilde{q}_{n,3/4}$ ainsi que
%la distance inter-quartile $\tilde{q}_{n,3/4} - \tilde{q}_{n,1/4}$.

\end{enumerate}

%\corrref{TD1_1}

\end{exercice}
